%% =========================================================
%% Chapter: Comparing marginal phases
%% Source: OpenAI, "A two-dimensional area law from a global spectral gap"
%%         (September 24, 2026), Lemma 5.2, 04-conditional.tex lines 310-479
%% =========================================================

\chapter{Comparing Marginal Phases}
\label{ch:marginal_phase}

This chapter formalizes the resolvent argument behind the marginal-phase
comparison of~\cite[Lemma~5.2]{OpenAI2026AreaLaw}: a change in products of
marginal imaginary powers is controlled by a conditional mutual information,
with constants depending only on the dimension of the moved system. The
faithful case is proved here; the passage to singular states is not yet
formalized.

\section{The resolvent defect of a compression}

Let $A$ and $B$ be positive definite operators with orthonormal eigenbases,
$V$ an isometry with $V^*AV=B$, $b_0$ a vector and $a_0=Vb_0$. For $v>0$ put
$h(v)=\langle a_0,(A+v)^{-1}a_0\rangle-\langle b_0,(B+v)^{-1}b_0\rangle$ and
$\delta_v=(A+v)^{-1}a_0-V(B+v)^{-1}b_0$.

\begin{lemma}[Resolvent compression identity]
    \label{lem:resolvent_compression}
    \lean{Matrix.ResolventCompression.discrepancy_quadratic,
        Matrix.ResolventCompression.defect_nonneg,
        Matrix.ResolventCompression.mul_norm_discrepancy_sq_le}
    \leanok
    $\langle\delta_v,(A+v)\delta_v\rangle=h(v)$. In particular $h(v)\ge0$ and
    $v\lVert\delta_v\rVert^2\le h(v)$.
\end{lemma}

\begin{proof}\leanok
    Expand the quadratic form and use $V^*(A+v)V=B+v$ and $V^*V=1$.
\end{proof}

\begin{lemma}[Tails and integral of the resolvent defect]
    \label{lem:resolvent_defect_integral}
    \lean{Matrix.ResolventCompression.defect_le_invQuadA,
        Matrix.ResolventCompression.defect_le_quadB_div_sq,
        Matrix.ResolventCompression.integrableOn_defect_and_integral}
    \leanok
    \uses{lem:resolvent_compression}
    For $v>0$, $h(v)\le\langle a_0,A^{-1}a_0\rangle$ and
    $h(v)\le\langle b_0,Bb_0\rangle/v^2$. Moreover $h$ is integrable on
    $(0,\infty)$ and
    $\int_0^\infty h(v)\,dv=\langle b_0,\log B\,b_0\rangle-\langle a_0,\log A\,a_0\rangle$.
\end{lemma}

\begin{proof}\leanok
    In eigencoordinates $h(v)=\sum_kp_k/(\lambda_k+v)-\sum_jq_j/(\mu_j+v)$ with
    $\sum_kp_k=\sum_jq_j$, and
    $\int_0^\infty\bigl((\lambda+v)^{-1}-(1+v)^{-1}\bigr)dv=-\log\lambda$.
\end{proof}

\begin{lemma}[Resolvent integral of a complex power]
    \label{lem:resolvent_mellin}
    \lean{Complex.integral_cpow_div_add, Complex.integral_cpow_imag_mul_inv_sub}
    \leanok
    For $\lambda>0$ and $-1<\operatorname{Re}z<0$,
    $\int_0^\infty v^z(\lambda+v)^{-1}dv=-\frac{\pi}{\sin\pi z}\lambda^z$. For
    real $u\ne0$,
    $\int_0^\infty v^{-iu}\bigl((\lambda+v)^{-1}-(1+v)^{-1}\bigr)dv
    =-\frac{\pi}{\sin(-i\pi u)}(\lambda^{-iu}-1)$.
\end{lemma}

\begin{proof}\leanok
    The substitution $v=\lambda t/(1-t)$ gives the Beta integral
    $B(z+1,-z)=\Gamma(z+1)\Gamma(-z)$, and the reflection formula evaluates it.
    The second identity follows from the first by dominated convergence as
    $\operatorname{Re}z\uparrow0$.
\end{proof}

\begin{theorem}[Phase gap of a compression]
    \label{thm:resolvent_phase_gap}
    \lean{Matrix.ResolventCompression.norm_phaseGap_le_integral,
        Matrix.ResolventCompression.norm_phaseGap_le_phaseRate}
    \leanok
    \uses{lem:resolvent_defect_integral, lem:resolvent_mellin}
    For every real $u$,
    $\lVert A^{-iu}a_0-VB^{-iu}b_0\rVert\le\frac{|\sinh\pi u|}{\pi}
    \int_0^\infty\sqrt{h(v)/v}\,dv$. If
    $\langle a_0,A^{-1}a_0\rangle\le D^2$, $\langle b_0,Bb_0\rangle\le1$ and
    $D\ge1$, then
    $\lVert A^{-iu}a_0-VB^{-iu}b_0\rVert\le2|\sinh\pi u|\,\mathcal R_\eta$ with
    $\eta=\int_0^\infty h$, $y=\min(1,\eta)$,
    $\mathcal R_\eta=\sqrt y+\sqrt{\eta\log(eD/y)}$ and $\mathcal R_0=0$.
\end{theorem}

\begin{proof}\leanok
    Pair the phase gap with a vector and use the second resolvent integral.
    Split $\int\sqrt{h/v}$ at $y/D^2$ and $1/y$; bound the ends by the tails
    and the middle by the arithmetic--geometric mean inequality.
\end{proof}

\section{The faithful case}

\begin{lemma}[Domination by the marginal]
    \label{lem:marginal_domination}
    \lean{Matrix.PosSemidef.le_card_smul_one_kronecker_partialTraceLeft}
    \leanok
    For positive semidefinite $\rho$ on $X\otimes Y$,
    $\rho\le(\dim X)\,(I_X\otimes\operatorname{tr}_X\rho)$.
\end{lemma}

\begin{proof}\leanok
    Write $\rho=B^*B$ and apply the Cauchy--Schwarz inequality over the
    $\dim X$ blocks.
\end{proof}

\begin{theorem}[Marginal-phase comparison, faithful case]
    \label{thm:marginal_phase_faithful}
    \lean{Entropy.MarginalPhase.norm_phaseDifference_le_of_posDef}
    \leanok
    \uses{thm:resolvent_phase_gap, lem:marginal_domination}
    Let $\rho$ be a positive definite density matrix on $xUF$, $d=\dim x$ and
    $\eta=I(x:F|U)_\rho$. For every real $u$,
    \[
      \bigl\lVert\bigl((\rho_U^{-iu}\rho_{xU}^{iu})\otimes I_F
        -\rho_{UF}^{-iu}\rho_{xUF}^{iu}\bigr)\sqrt\rho\bigr\rVert_2
      \le2|\sinh\pi u|\,\mathcal R_\eta .
    \]
\end{theorem}

\begin{proof}\leanok
    On Hilbert--Schmidt spaces take $A=L_{\sigma_f}R_{\rho^{-1}}$ and
    $B=L_{\sigma_s}R_{\rho_{xU}^{-1}}$ with $\sigma_f=I_x/d\otimes\rho_{UF}$ and
    $\sigma_s=I_x/d\otimes\rho_U$, the isometry
    $V(Y\sqrt{\rho_{xU}})=(Y\otimes I_F)\sqrt\rho$ and $b_0=\sqrt{\rho_{xU}}$.
    Then $\langle a_0,A^{-1}a_0\rangle=\operatorname{tr}\rho\sigma_f^{-1}\rho\le d^2$
    by the domination $\rho\le d^2\sigma_f$, $\langle b_0,Bb_0\rangle=1$, and
    the logarithmic gap is $I(x:F|U)_\rho$. The phase gap is the displayed
    vector up to a unimodular scalar.
\end{proof}
